% !PRINTED 211
\weberpage{211}
Jede Substitution $S$ läßt sich durch Zusammensetzung von

\[
\begin{pmatrix} m, & 0 \\ 0, & 1 \end{pmatrix}
\]

mit einer Substitution $T$ herleiten; man hat, damit

\[
S = \begin{pmatrix} m, & 0 \\ 0, & 1 \end{pmatrix} T
\]

sei, $\alpha, \beta, \gamma, \delta$ einfach aus den Kongruenzen

\[
a \equiv \alpha m, \quad b \equiv \beta m, \quad c \equiv \gamma, \quad d \equiv \delta \pmod{n}
\]

zu bestimmen.

Wir betrachten nun

\weq{10}{x_{\mu,\mu'} = \mathrm{sn}\left(\frac{4\mu K + 4\mu' i K'}{n}\right) = \frac{\vartheta_{00}\,\vartheta_{11}\left(\dfrac{2\mu + 2\mu'\omega}{n}\right)}{\vartheta_{10}\,\vartheta_{01}\left(\dfrac{2\mu + 2\mu'\omega}{n}\right)},}

\[
\varkappa^2 = \frac{\vartheta_{10}^4}{\vartheta_{00}^4}
\]

als Funktionen von $\omega$, und wenden darauf eine lineare Transformation

\[
\begin{pmatrix} \alpha, & \beta \\ \gamma, & \delta \end{pmatrix} \equiv \begin{pmatrix} 1, & 0 \\ 0, & 1 \end{pmatrix} \pmod{8}
\]

an, indem wir $\omega$ durch

\weq{11}{\omega' = \frac{\gamma + \delta\omega}{\alpha + \beta\omega}}

ersetzen; dann ergibt sich nach den Formeln (3), (5) und (6), \S\, 39, daß

\[
x_{\mu,\mu'} \quad \text{in} \quad x_{\alpha\mu + \gamma\mu',\,\beta\mu + \delta\mu'}
\]

übergeht, d.\,h. es erleiden die $x_{\mu,\mu'}$ eine Substitution, die nach (3), (6) mit

\[
T = \begin{pmatrix} \delta, & -\gamma \\ -\beta, & \alpha \end{pmatrix}
\]

zu bezeichnen ist, während $\varkappa^2$ ungeändert bleibt.

Auf diese Weise kann auch umgekehrt jede der Substitutionen $T$ entstehen; denn wenn irgend eine Substitution $T$ gegeben ist, so kann man durch Hinzufügen passender Vielfachen der (ungeraden) Zahl $n$ zu den Zahlen $\alpha, \beta, \gamma, \delta$ den Bedingungen

\[
\alpha \equiv 1, \quad \delta \equiv 1, \quad \beta \equiv 0, \quad \gamma \equiv 0 \pmod{8}
\]

immer genügen.
